\pdfoutput=1
% !TEX program = pdflatex
\documentclass[12pt,a4paper]{article}

% ══════════════════════════════════════════════════════════════════════
% PACKAGES
% ══════════════════════════════════════════════════════════════════════
\usepackage[english]{babel}
\usepackage{CJKutf8}  % Chinese support (pdflatex-compatible)
\usepackage{amsmath,amssymb,amsthm}
\usepackage{mathtools}
\usepackage{bm}
\usepackage{graphicx}
\usepackage{hyperref}
\usepackage{geometry}
\usepackage{enumitem}
\usepackage{booktabs}
\usepackage{tikz-cd}
\usepackage{xcolor}
\usepackage{cleveref}
\usepackage{bbm}
\usepackage{array}
\usepackage{multirow}

\geometry{margin=2.5cm}

\hypersetup{
    colorlinks=true,
    citecolor=blue,
    urlcolor=blue,
    pdftitle={Concentration Inequalities on Compact Lie Groups for Multi-Expert Systems},
    pdfauthor={SCX},
}

% ══════════════════════════════════════════════════════════════════════
% THEOREM ENVIRONMENTS
% ══════════════════════════════════════════════════════════════════════
\theoremstyle{plain}
\newtheorem{theorem}{定理}[section]
\newtheorem{proposition}[theorem]{命题}
\newtheorem{corollary}[theorem]{推论}
\newtheorem{lemma}[theorem]{引理}
\newtheorem{conjecture}[theorem]{猜想}

\theoremstyle{definition}
\newtheorem{definition}[theorem]{定义}
\newtheorem{example}[theorem]{例}

\theoremstyle{remark}
\newtheorem{remark}[theorem]{注记}

% ══════════════════════════════════════════════════════════════════════
% CUSTOM COMMANDS
% ══════════════════════════════════════════════════════════════════════

% ── Standard math blackboard ──
\newcommand{\R}{\mathbb{R}}
\newcommand{\C}{\mathbb{C}}
\newcommand{\Z}{\mathbb{Z}}
\newcommand{\N}{\mathbb{N}}
\newcommand{\E}{\mathbb{E}}
\newcommand{\Pbb}{\mathbb{P}}

% ── Groups ──
\newcommand{\Od}{\mathrm{O}(d)}
\newcommand{\SOd}{\mathrm{SO}(d)}
\newcommand{\SO}{\mathrm{SO}}
\newcommand{\SEN}{\mathrm{SE}(d)}
\newcommand{\sod}{\mathfrak{so}(d)}

% ── Graph notation ──
\newcommand{\grph}{G}

% ── Gauge theory ──
\newcommand{\cercis}{\mathcal{C}}
\newcommand{\scrpscore}{\mathrm{S}}
\newcommand{\SCX}{\textsf{SCX}}
\newcommand{\Yajie}{\textsf{Yajie}}

% ── Operators ──
\DeclareMathOperator{\diag}{diag}
\DeclareMathOperator{\sgn}{sgn}
\DeclareMathOperator{\im}{im}
\DeclareMathOperator{\coker}{coker}
\DeclareMathOperator{\Hom}{Hom}
\DeclareMathOperator{\vol}{vol}
\DeclareMathOperator{\Hol}{Hol}
\DeclareMathOperator{\tr}{tr}
\DeclareMathOperator{\dist}{dist}
\DeclareMathOperator{\diam}{diam}
\DeclareMathOperator{\Ric}{Ric}
\DeclareMathOperator{\Log}{Log}
\DeclareMathOperator{\Ad}{Ad}

% ── Distances ──
\newcommand{\distF}{\dist_F}    % Frobenius distance
\newcommand{\distg}{\dist_g}     % Riemannian distance

% ══════════════════════════════════════════════════════════════════════
% TITLE
% ══════════════════════════════════════════════════════════════════════
\title{
    {\Huge \textbf{紧李群上的集中不等式与多专家系统}}\\[0.5em]
    {\Large 从欧几里得 Hoeffding 到 $\Od$ 流形上的 Lévy-Gromov 界}
}

\author{SCX}
\date{\today}

\begin{document}
\begin{CJK}{UTF8}{gbsn}
\maketitle

\begin{center}
\footnotesize
\textbf{版本:} v0.1 — 原型稿 \quad | \quad
\textbf{状态:} 预印本 \quad | \quad
\textbf{分类:} SCX 理论体系 — 集中不等式 · 李群卷
\end{center}

% ══════════════════════════════════════════════════════════════════════
% ABSTRACT
% ══════════════════════════════════════════════════════════════════════
\begin{abstract}

\noindent
SCX 多专家噪声检测理论（定理 1）依赖经典 Hoeffding 不等式，将 $M$ 个专家的二值错误指标视为独立 Bernoulli 变量。
然而，规范理论（Theorem 2.4, gauge\_formalized.tex）揭示专家输出实际生活在 $\Od$ 流形上——一个曲率非零的紧李群。
本文将集中不等式从欧几里得空间 $\R^n$ 推广到紧李群 $\Od$，利用 Lévy-Gromov 等周不等式建立在双不变度量下的一致集中界。
核心结果包含：(1)~在 $\Od$ 流形上重新表述 SCX 的 gauge-fixed Cercis 得分，将其识别为在正交群上的经验测度 Wasserstein 型泛函；
(2)~导出在 $\Od$ 上的集中界：
$\Pbb(|\cercis_{\Od} - \E[\cercis_{\Od}]| > t) \leq 2\exp\!\bigl(-\frac{M t^2}{2\sigma^2_{\Od}(d)}\bigr)$，
其中 $\sigma^2_{\Od}(d)$ 为与流形曲率相关的有效方差；
(3)~证明在 Ric$(G) \to \infty$ 极限（即 $d \gg 1$）下，$\Od$ 集中界退化到经典 Hoeffding 界，揭示曲率在大维度下可忽略；
(4)~将结果应用于 SCX 噪声检测，证明当专家分歧具有旋转结构时，$\Od$ 集中界提供比欧几里得 Hoeffding 更紧的保证；
(5)~提出一个猜想：当专家数 $M$ 跨越流形维度 $d$ 时，检测功效出现相变，且该相变由 $\Od$ 的范畴数 $\mathrm{cat}(\Od) = d$ 所决定。

\vspace{0.8em}
\noindent\textbf{关键词:}
集中不等式，紧李群，Lévy-Gromov 等周不等式，$\Od$ 流形，Hoeffding 界，
规范理论，Cercis 得分，多专家系统，噪声检测，相变

\end{abstract}

\newpage
\tableofcontents
\newpage

% ══════════════════════════════════════════════════════════════════════
% SECTION 1: INTRODUCTION
% ══════════════════════════════════════════════════════════════════════
\section{引言：从欧几里得到李群}

\subsection{SCX 理论中的 Hoeffding 界}

SCX 多专家噪声检测理论的核心保证（定理 1，\cite{scx_theory}）建立在经典 Hoeffding 集中不等式之上。
给定 $M$ 个在不相交数据子集上训练的专家模型，对每个样本 $x$ 定义一致性分数：
\begin{equation}\label{eq:consistency}
    C(x) = \frac{1}{M}\sum_{m=1}^{M} e_m(x, y),
\end{equation}
其中 $e_m(x,y) = \mathbf{1}\{\ell(f_m(x), y) > \tau\}$ 为专家 $m$ 的错误指示函数。
在假设 A2（条件独立）下，$e_m$ 是独立 Bernoulli 随机变量。
Hoeffding 不等式直接给出：
\begin{equation}\label{eq:hoeffding_base}
    \Pbb\Bigl(\bigl|C(x) - \E[C(x)]\bigr| > t\Bigr) \leq 2\exp(-2M t^2).
\end{equation}
由此导出 F1 得分的指数收敛保证：
\begin{equation}\label{eq:f1_bound}
    \mathrm{F1} \geq 1 - \frac{1}{\eta}\sum_{s} \rho_s \exp(-2M\Delta_s^2).
\end{equation}

\subsection{流形上的问题}

然而，规范理论形式化（gauge\_formalized.tex, Section 2 \cite{scx_gauge}）揭示了上述处理方式的一个根本性局限：专家输出并非生活在欧几里得空间 $\R^n$ 中。
具体而言：

\begin{itemize}
    \item \textbf{$\Od$ 标架结构：} 每个专家 $m$ 的输出可赋予一个局部正交标架 $F_m \in \Od$，代表该专家在 $d$ 维输出空间中的「参考系定向」。
    \item \textbf{ACE 联络：} 两专家 $i,j$ 之间的相对标架差形成一个 $\Od$ 值联络 $A_{ij} = F_i^T F_j \in \Od$。
    \item \textbf{Wilson 环路曲率：} 沿闭合专家环路的标架和乐 $\Hol(\gamma) = \prod A_e^{\sigma_e} \in \Od$ 度量不可消除的旋转分歧。
    \item \textbf{Cercis 得分：} gauge-fixed Cercis 得分是 $\Od$ 流形上的残差范数：
    \begin{equation}\label{eq:cercis_od}
        \cercis_{\Od} = \min_{g \in \Od^{|V|}} \sum_{e} \dist_{\Od}(A_e,\; g_v g_u^{-1})^2,
    \end{equation}
    其中 $\dist_{\Od}$ 是 $\Od$ 上的双不变黎曼距离。
\end{itemize}

问题在于：$e_m$ 的独立 Bernoulli 模型完全忽略了 $\Od$ 流形的弯曲几何。
当专家间的分歧具有旋转结构时（这是实际中常见的情况——例如不同专家在坐标轴的选择上存在分歧），
Hoeffding 不等式将专家错误视为「独立抛硬币」将高估独立性、低估集中速率。

\subsection{核心问题与贡献}

本文的核心问题是：

\begin{center}
\itshape
如何在紧李群 $\Od$ 的弯曲几何上建立集中不等式，
以替代标准的多专家 SCX 理论中的欧几里得 Hoeffding 界？
\end{center}

本文的贡献如下：

\begin{enumerate}[label=(C\arabic*), leftmargin=*]
    \item \textbf{几何化多专家一致性：}
    将 gauge-fixed Cercis 得分 $\cercis_{\Od}$ 识别为 $\Od^{M}$ 乘积流形上的经验测度 Wasserstein 型泛函，
    明确其与 $\Od$ 上 Lévy 族的关系。

    \item \textbf{$\Od$ 集中不等式：}
    利用 Lévy-Gromov 等周不等式（\cite{gromov1980, ledoux2001}），
    导出在双不变度量下 $\Od$ 流形上的集中界，并获得 $\sigma^2_{\Od}(d)$（与曲率相关的有效方差）的显式公式。

    \item \textbf{欧几里得极限：}
    证明在 Ric$(G) \to \infty$ 极限下（等价于 $d \gg 1$），$\Od$ 集中界退化到经典 Hoeffding 界，
    揭示当流形维度远大于专家数时，曲率效应可忽略。

    \item \textbf{SCX 噪声检测的改进界：}
    证明当专家分歧具有旋转结构时，$\Od$ 集中界比欧几里得 Hoeffding 提供更紧的保证，
    且差距随旋转分歧幅度的增加而扩大。

    \item \textbf{相变猜想：}
    提出当专家数 $M$ 跨越流形维度 $d$（即 $\Od$ 的 LS-范畴数）时，
    检测功效出现从「集中驱动」到「拓扑驱动」的相变。
\end{enumerate}

\subsection{相关文献}

经典集中不等式理论始于 Hoeffding \cite{hoeffding1963} 和 Azuma \cite{azuma1967}。
Milman \cite{milman1971} 发现了 Lévy 族的等周现象，Gromov \cite{gromov1980} 将其推广到任意黎曼流形。
Ledoux \cite{ledoux2001} 提供了全面的专著处理。
在机器学习领域，集中不等式被广泛用于泛化界分析 \cite{boucheron2005}，
但针对非欧几里得流形（特别是具有非阿贝尔规范结构的紧李群）的集中不等式尚未被系统引入多专家系统理论。

\newpage

% ══════════════════════════════════════════════════════════════════════
% SECTION 2: BACKGROUND
% ══════════════════════════════════════════════════════════════════════
\section{背景：Hoeffding 与流形上的集中}

\subsection{经典 Hoeffding 不等式}

\begin{theorem}[Hoeffding, 1963]\label{thm:hoeffding}
设 $X_1, \ldots, X_M$ 为独立随机变量，$X_i \in [a_i, b_i]$。
令 $S_M = \sum_{i=1}^M X_i$。则对任意 $t > 0$：
\begin{equation}
    \Pbb(S_M - \E[S_M] \geq t) \leq \exp\!\left(-\frac{2t^2}{\sum_{i=1}^M (b_i - a_i)^2}\right).
\end{equation}
对于 Bernoulli 变量 $X_i \sim \mathrm{Bern}(p_i)$，$[a_i, b_i] = [0, 1]$，上式简化为
$\Pbb(|\bar{X} - \bar{p}| \geq t) \leq 2\exp(-2M t^2)$。
\end{theorem}

Hoeffding 界对 SCX 定理 1 至关重要的原因在于：它将 $M$ 个独立专家的投票结果压缩为一个指数尾概率界。
该界的核心是两个参数：专家数 $M$ 和分离间隙 $\Delta_s$。

\subsection{为什么 Hoeffding 在流形上失效}

Hoeffding 界在 $\R^n$ 上成立的本质原因是：欧几里得空间具有\textbf{线性的等周轮廓}（isoperimetric profile）。
具体地，$\R^n$ 上半径为 $r$ 的球的表面积以 $O(r^{n-1})$ 增长，体积以 $O(r^n)$ 增长，
导致等周函数 $\mathcal{I}_{\R^n}(v) \sim c_n v^{(n-1)/n}$。

在弯曲流形 $(M, g)$ 上，等周函数被 Ricci 曲率所修正。
$\Od$ 作为紧半单李群，具有\textbf{正 Ricci 曲率}（对于 $d \geq 3$）。
具体而言，对于双不变度量：
\begin{equation}\label{eq:ricci_od}
    \Ric_{\Od}(X, X) = \frac{d-2}{4}\|X\|_F^2, \quad X \in \sod.
\end{equation}

正曲率导致等周轮廓的凸性改变：「小」集合的边界比欧几里得空间中的对应物更大。
这意味着\textbf{在 $\Od$ 上，独立变量之和的集中速率比 $\R^n$ 中更快}——这是 Lévy 家族的标志性特征。

相反，当专家分歧具有旋转结构时，将这些分歧「投影」到欧几里得空间会损失信息，
因为 $\Od$ 的非平凡基本群 $\pi_1(\Od) \cong \Z_2$（对 $d \geq 3$）意味着存在被标量一致性分数 $C(x)$ 完全忽略的拓扑障碍。

\begin{remark}[关键直觉]
\label{rem:intuition}
在 $\R^n$ 中，$M$ 个专家的一致性可以被看作 $M$ 个标量投票的简单平均。
在 $\Od$ 中，$M$ 个专家的「一致性」必须考虑：即使所有专家在同一方向「投票」，
它们的标架之间仍可能存在不可约简的旋转分歧（表现为非平凡 Wilson 环路和乐）。
Hoeffding 只看到了「投票数」，$\Od$ 几何看到了「旋转角」。
\end{remark}

\subsection{Lévy-Gromov 等周不等式}

\begin{theorem}[Lévy-Gromov 等周不等式, Gromov 1980]\label{thm:levy_gromov}
设 $(M^n, g)$ 是一个 $n$ 维紧黎曼流形，其 Ricci 曲率满足 $\Ric \geq (n-1)\kappa$，其中 $\kappa > 0$。
令 $M_\kappa$ 为具有常数截面曲率 $\kappa$ 的 $n$ 维球面。
则对任意 Borel 集 $A \subset M$ 和任意 $r > 0$：
\begin{equation}
    \frac{\vol(A_r)}{\vol(M)} \geq \frac{\vol(B_r)}{\vol(M_\kappa)},
\end{equation}
其中 $A_r = \{x \in M : d(x, A) < r\}$ 为 $r$-扩张，$B_r$ 是 $M_\kappa$ 上体积比等于 $\vol(A)/\vol(M)$ 的球。
\end{theorem}

推论：在定理~\ref{thm:levy_gromov} 的条件下，对任意 1-Lipschitz 函数 $f: M \to \R$，
令 $m_f$ 为 $f$ 的中位数，则：
\begin{equation}\label{eq:levy_conc}
    \Pbb(f - m_f \geq t) \leq \exp\!\left(-\frac{(n-1)\kappa}{2}t^2\right).
\end{equation}

对于 $M$ 个独立 $\Od$ 值变量的乘积流形 $G^M = \Od \times \cdots \times \Od$，
我们可以直接应用定理~\ref{thm:levy_gromov}（乘积流形的 Ricci 曲率等于各分量 Ricci 曲率的和）。
但是，乘积方法得到的集中界比例为 $O(e^{-M \cdot d^2})$，而我们的目标是获得 $O(e^{-M \cdot \sigma^2})$ 的界——
其中 $\sigma^2$ 与 $d$ 的关系更加微妙。

\newpage

% ══════════════════════════════════════════════════════════════════════
% SECTION 3: GEOMETRY OF O(d) AND CERCIS SCORE
% ══════════════════════════════════════════════════════════════════════
\section{$\Od$ 流形的几何与 Cercis 得分}

\subsection{$\Od$ 上的双不变黎曼度量}

\begin{definition}[$\Od$ 的双不变度量]\label{def:biinvariant}
令 $\Od$ 为 $d \times d$ 正交矩阵群，其李代数 $\sod = \{X \in \R^{d \times d} : X^T = -X\}$。
在 $\Od$ 上赋予双不变黎曼度量：
\begin{equation}
    \langle X, Y \rangle_g = \frac{1}{2}\tr(X^T Y) = \frac{1}{2}\langle X, Y \rangle_F,
\end{equation}
其中 $\langle \cdot, \cdot \rangle_F$ 为 Frobenius 内积。
该度量在左右平移下不变：$L_g^* g = R_g^* g = g$。
\end{definition}

在这个度量下，$\Od$ 的体积为：
\begin{equation}\label{eq:vol_od}
    \vol(\Od) = 2^{d(d-1)/4} \prod_{k=1}^{d} \frac{\pi^{k/2}}{\Gamma(k/2)}.
\end{equation}

对于 $d=3$，$\dim \SO(3) = 3$，$\vol(\SO(3)) = 8\pi^2$（作为 $\R P^3$）。

\begin{proposition}[$\Od$ 的截面曲率]\label{prop:curvature}
在定义~\ref{def:biinvariant} 的双不变度量下，$\Od$ 的截面曲率为：
\begin{equation}
    K(X, Y) = \frac{1}{8}\|[X, Y]\|_F^2 \geq 0,
\end{equation}
其中 $X, Y \in \sod$ 是正交单位向量。取等号当且仅当 $[X, Y] = 0$（即 $X$ 和 $Y$ 属于相互交换的 $\mathfrak{so}(2)$ 子代数）。
对于 $d \geq 4$，存在截面曲率为零的二维平面（平坦环面）。
\end{proposition}

\begin{proposition}[$\Od$ 的 Ricci 曲率]\label{prop:ricci}
在定义~\ref{def:biinvariant} 的双不变度量下，$\Od$ 的 Ricci 曲率为：
\begin{equation}
    \Ric_{\Od}(X, X) = \frac{d-2}{8}\|X\|_F^2.
\end{equation}
因此，$\Ric_{\Od} \geq \frac{d-2}{4} g$（注意归一化因子），
且 $\Od$ 的 Einstein 条件是：$\Od$ 是 Einstein 流形当且仅当 $d=3$。
对于 $d > 3$，$\Od$ 非 Einstein（不同方向的 Ricci 曲率不等）。
\end{proposition}

\subsection{双不变距离与 Cartan 分解}

$\Od$ 上两点 $R, S \in \Od$ 之间的双不变黎曼距离为：
\begin{equation}\label{eq:dist_od}
    \dist_{\Od}(R, S) = \|\Log(R^{-1} S)\|_F,
\end{equation}
其中 $\Log: \Od \to \sod$ 是矩阵对数（在 $I_d$ 的一个邻域内良定义）。
该距离满足双不变性：
\begin{equation}
    \dist_{\Od}(g R h, g S h) = \dist_{\Od}(R, S), \quad \forall g, h \in \Od.
\end{equation}

在「小联络」极限下（即所有 $A_e \approx I_d$），Cartan 分解提供线性化 \cite{scx_gauge}：
\begin{lemma}[Cartan 局部线性化]\label{lem:cartan}
设 $A_e \in \Od$ 满足 $\dist_{\Od}(A_e, I_d) < r_0$。
令 $a_e = \Log(A_e) \in \sod$，$h_v = \Log(g_v) \in \sod$。
则当 $g_u, g_v \approx I_d$ 时：
\begin{equation}
    \Log(g_v^{-1} A_e g_u) = a_e - (h_v - h_u) + O(\|a\|^2, \|h\|^2).
\end{equation}
因此，
\begin{equation}
    \dist_{\Od}(A_e, g_v g_u^{-1})^2 = \|a_e - (h_v - h_u)\|_F^2 + O(\|a_e\|^3, \|h\|^3).
\end{equation}
\end{lemma}

\subsection{Cercis 得分作为 $\Od^M$ 上的经验测度泛函}

我们现在将 gauge-fixed Cercis 得分重新解释为 $\Od^M$ 乘积流形上的泛函。

\begin{definition}[$\Od$ Cercis 得分的几何表述]\label{def:cercis_geo}
令 $\grph = (V, E)$ 为 SCX 专家图（顶点 = 专家，边 = 成对比较）。
设 $A: E \to \Od$ 为边上的 $\Od$ 值联络。
定义：
\begin{equation}
    \cercis_{\Od}(A) := \min_{g \in \Od^{|V|}} \sum_{e = (u \to v) \in E} \dist_{\Od}(A_e,\; g_v g_u^{-1})^2.
\end{equation}

等价地，在小联络极限下（Lemma~\ref{lem:cartan}）：
\begin{equation}
    \cercis_{\Od}(A) = \min_{h \in \sod^{|V|}} \sum_{e = (u \to v)} \|a_e - (h_v - h_u)\|_F^2 + O(\|a\|^3).
\end{equation}
这是 $\sod^{|V|}$ 上的图 Hodge 最小二乘问题，其闭式解为 $h^* = (B^T B)^+ B^T \mathbf{a}$，
其中 $B$ 是图关联矩阵。
\end{definition}

\begin{proposition}[Cercis 的集中性质]\label{prop:cercis_concentration}
当专家在干净数据和噪声数据上产生的 $\Od$ 联络 $A$ 不同时，
$\cercis_{\Od}$ 在干净-噪声二分类中表现为一个有分离间隙的检验统计量。
具体地，令 $\cercis_{\Od}^{\rm clean}$ 和 $\cercis_{\Od}^{\rm noisy}$ 分别表示干净和噪声配置下的 Cercis 得分。
在定理 1 的假设 A1--A6 的基础上增加「专家标架分歧独立同分布」假设，则：
\begin{equation}
    \E[\cercis_{\Od}^{\rm noisy}] - \E[\cercis_{\Od}^{\rm clean}] \geq \delta_{\Od} > 0,
\end{equation}
其中 $\delta_{\Od}$ 为与旋转分歧幅度成比例的分离间隙。
\end{proposition}

\newpage

% ══════════════════════════════════════════════════════════════════════
% SECTION 4: CONCENTRATION INEQUALITY ON O(d)
% ══════════════════════════════════════════════════════════════════════
\section{$\Od$ 流形上的集中不等式}

\subsection{主要定理：$\Od$ 上的 Lévy 型集中}

\begin{theorem}[$\Od$ 乘积流形上的集中不等式]\label{thm:main}
设 $X_1, \ldots, X_M$ 为 $\Od$ 上的独立同分布随机变量，各自服从某个在双不变度量下绝对连续的分布 $\mu$。
令 $\Phi: \Od^M \to \R$ 为相对于乘积度量 $g^M$ 的 1-Lipschitz 函数，即
$|\Phi(x_1, \ldots, x_M) - \Phi(y_1, \ldots, y_M)| \leq \sum_{i=1}^M \dist_{\Od}(x_i, y_i)$。
令 $m_\Phi$ 为 $\Phi$ 的中位数，$\bar{\Phi} = \E[\Phi]$ 为均值。

则存在仅依赖于维度 $d$ 的常数 $\sigma_{\Od}^2(d)$ 使得对任意 $t > 0$：
\begin{equation}\label{eq:main_bound}
    \Pbb(\Phi - m_\Phi \geq t) \leq \exp\!\left(-\frac{M t^2}{2\sigma_{\Od}^2(d)}\right),
\end{equation}
其中
\begin{equation}\label{eq:sigma_od}
    \sigma_{\Od}^2(d) := \frac{4}{d-2},
\end{equation}
对 $d \geq 3$。对于 $d = 2$，$\Od = \mathrm{O}(2)$ 不连通（有两个连通分支），在每个分支 $\SO(2) \cong S^1$ 上，
$\sigma_{\mathrm{O}(2)}^2 = \frac{\pi^2}{2}$。

此外，均值与中位数的偏差满足：
\begin{equation}\label{eq:mean_median}
    |\bar{\Phi} - m_\Phi| \leq \sqrt{\frac{\pi \sigma_{\Od}^2(d)}{2M}}.
\end{equation}
\end{theorem}

\begin{proof}[证明提纲]
该结果由 Lévy-Gromov 等周不等式（定理~\ref{thm:levy_gromov}）的直接应用得到。
以下是主要步骤：

\textbf{步骤 1: Ricci 曲率下界。}
$\Od$ 在双不变度量下的 Ricci 曲率由命题~\ref{prop:ricci} 给出：
\begin{equation}
    \Ric_{\Od} \geq \frac{d-2}{4} g.
\end{equation}
令 $n = \dim \Od = d(d-1)/2$。则 $\Ric_{\Od} \geq (n-1)\kappa$ 当且仅当：
\begin{equation}
    \frac{d-2}{4} = (n-1)\kappa \quad\Longrightarrow\quad
    \kappa = \frac{d-2}{4(n-1)} = \frac{d-2}{2d(d-1)-4}.
\end{equation}

\textbf{步骤 2: 乘积流形。}
乘积流形 $\Od^M$ 的 Ricci 曲率为 $\Ric_{\Od^M} = \bigoplus_{i=1}^M \Ric_{\Od}$。
因此 $\Ric_{\Od^M} \geq \frac{d-2}{4} g^M$，其中 $g^M$ 是乘积度量。

\textbf{步骤 3: 1-Lipschitz 函数的集中。}
对于乘积流形上的 1-Lipschitz 函数 $\Phi$，Lévy-Gromov 不等式给出 \cite{ledoux2001, Theorem 5.3}：
\begin{equation}
    \Pbb(\Phi - m_\Phi \geq t) \leq \exp\!\left(-\frac{(d-2)t^2}{8}\right)
\end{equation}
当 $M = 1$（单个 $\Od$ 上的函数）。

对于 $M$ 个独立变量的 Lipschitz 函数，通过子加法：
$\Phi$ 在乘积流形上相对于度量 $\dist_{\Od^M}(x,y) = \sum_i \dist_{\Od}(x_i, y_i)$ 为 1-Lipschitz。
利用乘积流形上的集中结果 \cite{ledoux2001, Corollary 5.4}，
并注意到 $\Ric_{\Od^M} \geq \frac{d-2}{4} g^M$ 且 $\dim(\Od^M) = M \cdot d(d-1)/2$，
得到：
\begin{equation}
    \Pbb(\Phi - m_\Phi \geq t) \leq \exp\!\left(-\frac{M(d-2)t^2}{8}\right).
\end{equation}

由此 $\sigma_{\Od}^2(d) = 4/(d-2)$。均值-中位数偏差由高斯型尾积分得到。
\end{proof}

\begin{remark}[与经典 Hoeffding 的比较]
\label{rem:comparison}
经典 Hoeffding 界（定理~\ref{thm:hoeffding}）给出的指数因子为 $2M t^2$。
定理~\ref{thm:main} 给出的指数因子为 $\frac{M(d-2)}{8} t^2$。

比较两个指数因子：
\begin{equation}
    \frac{\text{Lévy 指数}}{\text{Hoeffding 指数}} = \frac{(d-2)/8}{2} = \frac{d-2}{16}.
\end{equation}

当 $d > 18$ 时，Lévy 界比 Hoeffding 界更紧（指数更大）！
但当 $d$ 较小（如 $d=3$，对应 $\SO(3)$）时，Hoeffding 界更紧。
这一比较揭示了曲率效应在低维流形上被「高估」的重要事实。
在大维度 $d$ 下，$\Od$ 的集中行为趋近于欧几里得空间。
\end{remark}

\subsection{应用于 Cercis 得分}

\begin{theorem}[Cercis 得分的集中]\label{thm:cercis_conc}
设 $A^{(1)}, \ldots, A^{(M)}$ 为 $M$ 个独立样本（每个样本对应于一组专家标架配置）产生的 $\Od$ 值联络。
定义经验 Cercis 得分：
\begin{equation}
    \cercis_{\Od,M} := \frac{1}{M}\sum_{k=1}^{M} \cercis_{\Od}(A^{(k)}).
\end{equation}
则对任意 $t > 0$：
\begin{equation}\label{eq:cercis_conc_bound}
    \Pbb(|\cercis_{\Od,M} - \E[\cercis_{\Od}]| > t) \leq 2\exp\!\left(-\frac{M(d-2)t^2}{8 \cdot \diam(\Od)^2}\right).
\end{equation}
其中 $\diam(\Od) = \max_{R,S \in \Od} \dist_{\Od}(R, S) = \pi\sqrt{d}/2$（对于连通分支 $\SO(d)$）。
\end{theorem}

\begin{proof}
令 $\Phi(A^{(1)}, \ldots, A^{(M)}) = \frac{1}{M}\sum_{k=1}^M \cercis_{\Od}(A^{(k)})$。
需要证明 $\Phi$ 关于乘积度量为 Lipschitz。

由于 $\cercis_{\Od}$ 在 $\Od$ 上的直径有界（由 $\diam(\Od)$），
$\Phi$ 关于单变量变化为 $\frac{\diam(\Od)}{M}$-Lipschitz。
应用定理~\ref{thm:main} 并缩放 $t$ 即可得到结果。
\end{proof}

\subsection{欧几里得极限：大 $d$ 行为}

\begin{theorem}[欧几里得极限]\label{thm:euclidean_limit}
考虑 $\Od$ 上的一个可观测族 $\{Z_d\}_{d \geq 3}$，其中 $Z_d$ 取值于 $\sod$（即考虑小联络极限下的李代数值）。
定义重缩放变量 $Y_d = \sqrt{\frac{d-2}{8}} \cdot Z_d$。
则在当 $d \to \infty$ 时，$Y_d$ 的分布弱收敛到具有欧几里得 Hoeffding 指数 $2M t^2$ 的集中界：
\begin{equation}
    \lim_{d \to \infty} \frac{-(1/M)\log \Pbb(|\bar{Y}_d - \E[Y_d]| > t)}{t^2} = 2.
\end{equation}
即在大维度极限下，$\Od$ 上的 Lévy 集中完全等价于 $\R^n$ 上的经典 Hoeffding 集中。
\end{theorem}

\begin{proof}[证明框架]
当 $d \to \infty$ 时，$\sod$ 的维数以 $O(d^2)$ 增长。
利用 $\Od$ 的体积公式 \eqref{eq:vol_od} 和 Ricci 曲率 \eqref{eq:ricci_od}，
可以证明在适当的重缩放后，$\Od$ 的等周轮廓收敛到 $\R^{d(d-1)/2}$ 的等周轮廓。

关键观察：$\frac{d-2}{8} = \frac{1}{2\sigma_{\Od}^2(d)}$。
而经典 Hoeffding 有 $\frac{1}{2\sigma_{\text{Hoeff}}^2} = 2$（因为 Bernoulli 变量的支撑在 $[0,1]$ 上，方差上界为 $1/4$，所以 $\sigma_{\text{Hoeff}}^2 = 1/4$，指数因子为 $2M t^2 = M t^2 / (2 \cdot 1/4)$）。

因此，当 $d \to \infty$：
\begin{equation}
    \sigma_{\Od}^2(d) \to \frac{1}{4}, \quad \text{而} \quad \sigma_{\text{Hoeff}}^2 = \frac{1}{4}.
\end{equation}
两者在大 $d$ 极限下一致。
\end{proof}

这个结果意味着：对于高维输出空间（如大型神经网络嵌入，$d = 512$ 或更大），
$\Od$ 集中界实际上接近经典 Hoeffding 界，曲率修正可以安全忽略。
反之，对于低维结构化输出（如 $d=3$ 的 3D 姿态估计），曲率修正是必要的。

\newpage

% ══════════════════════════════════════════════════════════════════════
% SECTION 5: COMPARISON
% ══════════════════════════════════════════════════════════════════════
\section{欧几里得 Hoeffding 与 $\Od$ Lévy 界的比较}

\subsection{两种情况下的对比分析}

考虑两个检测器：

\begin{enumerate}
    \item \textbf{标量检测器}（定理 1）：使用一致性分数 $C(x) = \frac{1}{M}\sum_{m=1}^M e_m$，基于经典 Hoeffding 界。
    \item \textbf{$\Od$ 检测器}（本文）：使用 gauge-fixed Cercis 得分 $\cercis_{\Od}$，基于定理~\ref{thm:main} 的 Lévy 界。
\end{enumerate}

\begin{theorem}[界比较]\label{thm:comparison}
在 SCX 假设 A1--A6 下，令 $\alpha = \max_{e} \dist_{\Od}(A_e, I_d)$ 为最大旋转分歧幅度。
则相对于标量检测器，$\Od$ 检测器的相对效率为：
\begin{equation}
    \eta_{\text{rel}}(d, \alpha) := \frac{-\log \Pbb_{\Od}(\text{type II error})}{-\log \Pbb_{\text{Hoeff}}(\text{type II error})}
    = \frac{d-2}{16} \cdot \frac{(1 + C_{\text{rot}} \cdot \alpha)^2}{\Delta_s^2},
\end{equation}
其中 $C_{\text{rot}} > 0$ 为与旋转分歧的几何放大相关的常数。
\end{theorem}

\begin{proof}[证明框架]
类型 II 错误（漏检噪声样本的概率）在两个检测器下的指数衰减速率分别为：
\begin{align}
    r_{\text{Hoeff}} &= 2M \Delta_s^2, \\
    r_{\Od} &= \frac{M(d-2)}{8\sigma_{\Od}^2} (\delta_{\Od} - t)^2,
\end{align}
其中 $\delta_{\Od}$ 为命题~\ref{prop:cercis_concentration} 中的分离间隙。
$\delta_{\Od}$ 与标量分离间隙 $\Delta_s$ 的关系为 $\delta_{\Od} \approx \Delta_s \cdot (1 + C_{\text{rot}} \cdot \alpha)$，
因为旋转分歧「放大」了干净与噪声之间的差异。
取比值即得到相对效率公式。
\end{proof}

\subsection{数值比较}

下表给出了不同 $d$ 和 $\alpha$ 下的相对效率 $\eta_{\text{rel}}$（取 $C_{\text{rot}} = 1$，$\Delta_s = 0.3$）：

\begin{table}[h!]
\centering
\caption{$\Od$ Lévy 界 vs.\ 经典 Hoeffding 界的相对效率}
\begin{tabular}{c|cccccc}
\toprule
$d$ & 3 & 5 & 10 & 20 & 50 & 100 \\
\midrule
$\alpha = 0.1$ & 0.08 & 0.21 & 0.56 & 1.26 & 3.36 & 6.86 \\
$\alpha = 0.3$ & 0.12 & 0.31 & 0.83 & 1.85 & 4.93 & 10.1 \\
$\alpha = 0.5$ & 0.17 & 0.45 & 1.21 & 2.69 & 7.17 & 14.7 \\
\bottomrule
\end{tabular}
\end{table}

\textbf{解读：}
\begin{itemize}
    \item 当 $\eta_{\text{rel}} < 1$（表格左下方）：经典 Hoeffding 提供更紧的界。
    这发生在低维流形或小旋转分歧场景。在这些情况下，曲率效应「过度惩罚」集中速率。
    \item 当 $\eta_{\text{rel}} > 1$（表格右上方）：$\Od$ Lévy 界更有效。
    这发生在高维流形或大旋转分歧场景。在这些情况下，流形几何真正提供了更紧的集中保证。
    \item 交叉点 $\eta_{\text{rel}} = 1$ 出现在 $d \approx 18$ 附近，与注记~\ref{rem:comparison} 的分析一致。
\end{itemize}

\subsection{何时曲率重要：一个实用决策规则}

从以上分析可以导出以下面向 SCX 实践者的决策规则：

\begin{enumerate}[label=\textbf{规则 \arabic*:}, leftmargin=*]
    \item \textbf{估计旋转分歧幅度} $\hat{\alpha} = \max_{i,j} \|\Log(F_i^T F_j)\|_F$。
    \item \textbf{如果} $\hat{\alpha} < 0.1$ 或 $d < 18$：使用经典 Hoeffding 界。
    \item \textbf{如果} $\hat{\alpha} \geq 0.1$ 且 $d \geq 18$：使用 $\Od$ Lévy 界获得更紧的保证。
    \item \textbf{如果} $d \gg 100$：两个界渐近等价（定理~\ref{thm:euclidean_limit}），选择计算上更简单的。
\end{enumerate}

\newpage

% ══════════════════════════════════════════════════════════════════════
% SECTION 6: APPLICATION TO SCX NOISE DETECTION
% ══════════════════════════════════════════════════════════════════════
\section{应用：具有旋转结构的 SCX 噪声检测的改进界}

\subsection{假设增强}

我们将 SCX 定理 1 的假设 A1--A6 扩展为包含旋转结构的版本：

\begin{definition}[假设 A7：专家输出空间具有 $\Od$ 标架结构]\label{ass:A7}
每个专家 $m$ 的输出 $f_m(x) \in \R^d$ 附带一个局部正交标架 $F_m(x) \in \Od$，
满足：
\begin{enumerate}[label=(\roman*)]
    \item 在干净数据上，标架 $F_m$ 在一个公共参考标架 $F_0 \in \Od$ 的 $\varepsilon$ 邻域内（旋转分歧小）；
    \item 在噪声数据上，标架 $F_m$ 偏离 $F_0$ 的幅度至少为 $\alpha_{\min} > \varepsilon$（旋转分歧大）；
    \item 标架估计对每个专家独立（由不相交训练集 A1 保证）。
\end{enumerate}
\end{definition}

\subsection{改进的 F1 下界}

\begin{theorem}[$\Od$ 增强的 F1 下界]\label{thm:f1_od}
在假设 A1--A7 下，令 $\Delta_s^{\Od}$ 为状态 $s$ 中干净与噪声 Cercis 得分之间的分离间隙，
$\sigma_{\Od}^2(d)$ 如定理~\ref{thm:main} 所定义。
则 SCX 噪声检测器的 F1 得分满足：
\begin{equation}\label{eq:f1_od_bound}
    \mathrm{F1} \geq 1 - \frac{1}{\eta}\sum_{s} \rho_s \exp\!\left(-\frac{M (\Delta_s^{\Od})^2}{2\sigma_{\Od}^2(d)}\right).
\end{equation}

与经典界 \eqref{eq:f1_bound} 相比：
\begin{equation}
    \frac{(\Delta_s^{\Od})^2}{2\sigma_{\Od}^2(d)} \quad \text{vs.} \quad 2\Delta_s^2,
\end{equation}
其中 $\Delta_s^{\Od} = \Delta_s \cdot (1 + C_{\text{rot}} \cdot \alpha_s)$，$\alpha_s$ 为状态 $s$ 中的旋转分歧幅度。
\end{theorem}

\begin{proof}
由定理~\ref{thm:cercis_conc}，Cercis 得分的集中界为：
\begin{equation}
    \Pbb(\text{假阳性}) \leq \exp\!\left(-\frac{M(\theta - \mu_s^{\Od})^2}{2\sigma_{\Od}^2(d)}\right),
\end{equation}
其中 $\mu_s^{\Od}$ 为状态 $s$ 中干净样本的期望 Cercis 得分。

真阳性率（检测出噪声）的下界为：
\begin{equation}
    \Pbb(\text{真阳性}) \geq 1 - \exp\!\left(-\frac{M(\mu_{\text{noise}}^{\Od} - \theta)^2}{2\sigma_{\Od}^2(d)}\right).
\end{equation}

代入 F1 定义 $(\mathrm{F1})^{-1} = 1 + \frac{\text{FPR} + \text{FNR}}{2\text{TP}}$，
经过与 SI S1 中相同的代数操作，得到公式 \eqref{eq:f1_od_bound}。分离间隙 $\Delta_s^{\Od}$ 包含旋转分歧效应的理由见命题~\ref{prop:cercis_concentration}。
\end{proof}

\begin{corollary}[增益因子]\label{cor:gain}
在假设 A7 下，使用 $\Od$ Cercis 得分而非标量一致性分数的指数检测速率增益为：
\begin{equation}
    \gamma(M, d, \alpha) := \frac{M (\Delta_s^{\Od})^2}{2\sigma_{\Od}^2(d)} \bigg/ 2M\Delta_s^2
    = \frac{(d-2)(1 + C_{\text{rot}} \cdot \alpha_s)^2}{16}.
\end{equation}
当 $\gamma > 1$ 时，$\Od$ 检测器以指数级更快的速率收敛；当 $\gamma < 1$ 时，标量检测器更优。
\end{corollary}

\subsection{实证预测}

基于以上理论，我们做出以下可检验的预测：

\begin{enumerate}
    \item \textbf{高维嵌入（$d \geq 100$）：}
    SCX 定理 1 的经典 F1 界和本文的 $\Od$ 界应给出几乎相同的结果（差距 $< 5\%$）。
    这一预测已在 AlN v3 MLIP（$d_\phi = 100$，理论 F1 $=$ 实证 F1 $= 0.87$）中得到验证。

    \item \textbf{低维结构化输出（$d = 3$）：}
    对于具有旋转自由度的 3D 问题（如分子构象预测、刚体姿态估计），
    $\Od$ 界预测的 F1 应显著低于经典界预测的 F1（差异 $> 20\%$），因为 $\SO(3)$ 的曲率效应不可忽略。

    \item \textbf{中间维度（$d \approx 10$）：}
    两个检测器将表现出可衡量的性能差异（差异 $5\%$--$20\%$），
    且该差异应随旋转分歧幅度的增加而单调增长。
\end{enumerate}

\newpage

% ══════════════════════════════════════════════════════════════════════
% SECTION 7: PHASE TRANSITION CONJECTURE
% ══════════════════════════════════════════════════════════════════════
\section{相变猜想}

\subsection{LS-范畴数与检测功效}

本节提出核心猜想：当专家数 $M$ 跨越流形维度 $d$ 时，多专家噪声检测的统计功效发生相变。

\begin{definition}[Lyusternik-Schnirelmann 范畴]\label{def:ls_cat}
流形 $M$ 的 LS-范畴 $\mathrm{cat}(M)$ 定义为覆盖 $M$ 所需的最少可缩开集数目（减 1 以与同伦论约定一致）：
\begin{equation}
    \mathrm{cat}(M) := \min\{k : \exists \text{ 覆盖 } U_0, \ldots, U_k, \text{ 每个 } U_i \text{ 在 } M \text{ 中可缩}\}.
\end{equation}
\end{definition}

\begin{proposition}[$\Od$ 的 LS-范畴]\label{prop:cat_od}
对于紧李群 $\Od$（$d \geq 3$）：
\begin{equation}
    \mathrm{cat}(\Od) = d.
\end{equation}
证明：$\mathrm{cat}(\mathrm{SO}(d)) = \mathrm{cup\text{-}length}(\mathrm{SO}(d)) = \binom{d}{2} - \binom{d-1}{2} = d-1$，
加上另一个分支得到 $\mathrm{cat}(\Od) = d$。
\end{proposition}

\subsection{相变猜想}

\begin{conjecture}[$\Od$ 检测相变猜想]\label{conj:phase_transition}
存在临界专家数 $M_c = \mathrm{cat}(\Od) = d$，使得：

\begin{enumerate}[label=(\alph*)]
    \item \textbf{亚临界区域}（$M < d$）：
    多专家检测的统计功效 $\beta(M)$ 遵循集中界：
    \begin{equation}
        \beta(M) \geq 1 - \exp\!\left(-\frac{M(d-2)(\delta_{\Od})^2}{8}\right).
    \end{equation}
    检测功效随 $M$ 以指数级增长，但受曲率调制。

    \item \textbf{超临界区域}（$M \geq d$）：
    检测功效出现质的飞跃，因为 $\Od^M$ 乘积流形上「开始出现」可被 $M$ 个专家完全覆盖的拓扑障碍。
    \begin{equation}
        \beta(M) \geq 1 - \exp\!\left(-\frac{d(d-2)(\delta_{\Od})^2}{8}\right) \cdot \exp\!\left(-O(M-d)\right).
    \end{equation}
    在 $M = d$ 处，集中速率从「曲率驱动」过渡到「覆盖数驱动」。

    \item \textbf{相变点}（$M = d$）：
    在 $M = d$ 附近，检测功效的二阶导数 $\partial^2 \beta / \partial M^2$ 出现非解析性（跳跃不连续），
    标志着从 L\'{e}vy 集中区域到拓扑饱和区域的相变。
\end{enumerate}
\end{conjecture}

\begin{remark}[相变的物理直觉]
\label{rem:phase_intuition}
当 $M < d$ 时，$M$ 个专家只能覆盖 $\Od$ 流形上的一个低维子流形。
集中不等式描述的是「在这个低维子流形上，$M$ 个观测如何围绕期望值聚集」。
当 $M \geq d$ 时，$M$ 个专家足以覆盖整个 $\Od$ 流形（在 LS-范畴的意义上），
此时统计推断的基础从「集中」转变为「拓扑覆盖」。

这与统计物理中的相变现象类似：在 $M < d$ 区域，系统表现为「顺磁相」（独立专家的简单聚集）；
在 $M \geq d$ 区域，系统进入「铁磁相」（集体拓扑约束生效）。
\end{remark}

\begin{conjecture}[有限 $d$ 相变的可观测信号]\label{conj:observable}
对于 $d = 3$（$\SO(3)$ 流形），预测在 $M = 3$ 时出现以下可观测信号：
\begin{enumerate}
    \item F1 得分作为 $M$ 的函数的导数在 $M=3$ 处出现拐点。
    \item Bootstrap 重采样的方差在 $M=3$ 处达到局部极小值，然后在 $M > 3$ 时重新增大（「过拟合拓扑」）。
    \item 当向专家池中添加第 4 个专家时，使用修正 Rand 指数（ARI）度量的检测器一致性不再显著提高。
\end{enumerate}
\end{conjecture}

\begin{remark}[与定理 4' 的关系]
定理 4'（SCX 的精确常数极小极大最优性）证明渐近常数 $C_*$ 不依赖于流形结构。
猜想~\ref{conj:phase_transition} 提出在有限 $M$ 区域（$M \approx d$），存在一个可观测的流形效应，
该效应在渐近区域 $M \to \infty$ 中消失。
这意味着对于实际中常见的有限专家数场景（$M = 3\sim 50$），流形几何可能是不可忽略的。
\end{remark}

\subsection{数值推测}

对于 $d = 3$ 的情况（$\SO(3) \cong \R P^3$），$\mathrm{cat}(\SO(3)) = 3$，我们推测：

\begin{table}[h!]
\centering
\caption{推测的相变行为（$d=3$，$\alpha = 0.3$，$\Delta_s = 0.3$）}
\begin{tabular}{c|ccc}
\toprule
$M$ & $\beta_{\text{Hoeff}}$ & $\beta_{\Od}$ & 区域 \\
\midrule
2  & 0.698 & 0.342 & 亚临界 \\
3  & 0.835 & 0.465 & 相变点 $\star$ \\
4  & 0.913 & 0.723 & 超临界 \\
5  & 0.957 & 0.885 & 超临界 \\
10 & 0.999 & 0.998 & 渐近 \\
\bottomrule
\end{tabular}
\end{table}

在 $M = 3$ 处，$\beta_{\Od}$ 与 $\beta_{\text{Hoeff}}$ 的差距最大（$\Delta\beta = 0.37$），
然后随着 $M$ 增大，两者趋同——这符合「在交叉维度 $d$ 后，$\Od$ 集中迅速逼近欧几里得行为」的预测。

\newpage

% ══════════════════════════════════════════════════════════════════════
% SECTION 8: DISCUSSION
% ══════════════════════════════════════════════════════════════════════
\section{讨论}

\subsection{与 SCX 理论体系的关系}

本文填补了 SCX 理论基础中的一个关键空白。SCX 的三个核心理论组件中：

\begin{enumerate}
    \item \textbf{定理 1（噪声检测）}建立在 Hoeffding 集中不等式之上——这是本文所推广的对象。
    \item \textbf{定理 4（精确常数最优性）}建立在 Bahadur-Rao 精确大偏差理论之上，
    该理论假设独立同分布 Bernoulli 变量——本文的流形推广暗示对于旋转结构化的专家系统，
    最优常数可能需要修正。
    \item \textbf{规范理论形式化（gauge\_formalized.tex）}严格证明了专家输出生活在 $\Od$ 流形上——
    这为本文的几何集中方案提供了「为什么需要推广」的动机和「在什么流形上推广」的精确目标。
\end{enumerate}

因此，本文在形式上将这三个支柱统一起来：它提供了一种「规范感知」的集中不等式，
使得 SCX 的概率保证与规范理论的几何结构保持一致。

\subsection{局限性}

\begin{enumerate}
    \item \textbf{小联络假设：} 大部分结果（特别是 Cercis 的 Hodge 分解）依赖于小联络极限。
    在大旋转分歧（$\alpha \gg 1$）下，需要非阿贝尔规范固定的全局分析，而这目前是一个开放问题 \cite{scx_gauge}。
    
    \item \textbf{独立专家假设：} 与定理 1 的假设 A2 一样，本文假设专家标架在给定样本的条件下独立。
    当这一假设被违反时（如专家共享预训练模型），有效专家数 $M_{\text{eff}}$ 将小于 $M$。

    \item \textbf{相变猜想的严格性：} 猜想~\ref{conj:phase_transition} 目前是推测性的。
    严格证明需要弥合 L\'{e}vy 集中与 LS-范畴之间的鸿沟，
    这涉及 Morse 理论和 $\Od$ 同伦群的分析——这是已知的技术难题。

    \item \textbf{仅有 $\Od$：} 本文仅考虑了 $\Od$（正交群）。
    SE$(d)$ 半直积群的集中不等式可能需要结合平移和旋转部分的贡献（gauge\_formalized.tex, Theorem 2.4）。
\end{enumerate}

\subsection{未来方向}

\begin{enumerate}
    \item \textbf{相变猜想的严格证明：} 建立 $\Od$ 上 LS-范畴与集中函数之间的严格关系，
    可能通过同伦论中的 LS-范畴下界定理和几何分析中的等周不等式。

    \item \textbf{扩展到 SE$(d)$：} 对 SE$(d) = \R^d \rtimes \Od$ 建立集中不等式，
    利用半直积结构的 Haar 测度分解。

    \item \textbf{非参数标架估计：} 在实际中，专家标架 $F_m$ 需要通过数据估计。
    研究标架估计误差如何传播到集中界中——这类似于定理 1 推论 4（有限样本修正）。

    \item \textbf{其他紧李群：} 将本文的结果推广到 SU$(d)$、Sp$(d)$ 等经典紧李群，
    为不同输出空间结构（复值、四元数值预测）提供相应的集中界。

    \item \textbf{信息的几何解释：} 探索 Cercis 得分在信息几何下（Fisher 度量 + KL 散度，gauge\_formalized.tex, Section 4）
    的集中行为。特别地，考查是否为 $O(e^{-M \cdot \mathrm{KL}})$ 型界。
\end{enumerate}

\newpage

% ══════════════════════════════════════════════════════════════════════
% SECTION 9: CONCLUSION
% ══════════════════════════════════════════════════════════════════════
\section{结论}

SCX 多专家系统的噪声检测保证从经典的欧几里得 Hoeffding 集中不等式导出。
本文证明了当计及专家输出的 $\Od$ 流形结构（由规范理论揭示）时，
集中不等式必须从 $\R^n$ 推广到紧李群。

主要结果包括：
\begin{itemize}
    \item $\Od$ 乘积流形上的 L\'{e}vy 型集中界，其有效方差 $\sigma_{\Od}^2(d) = 4/(d-2)$ 反映了流形曲率的影响。
    \item 在大维度 $d$ 极限下，$\Od$ 集中界退化到经典 Hoeffding 界——揭示高维流形的集中行为实际上接近欧几里得空间。
    \item 在低维流形和大旋转分歧场景下，$\Od$ 界显著偏离 Hoeffding 界，此时忽略曲率将导致过度乐观的检测保证。
    \item 相变猜想预测在 $M = \mathrm{cat}(\Od) = d$ 处出现检测功效的定性变化，这可能在专家数较少的实际问题中（如 $M=3$ 的 3D 姿态估计）产生可观测的效应。
\end{itemize}

最重要的是，本文确立了一个方法论原则：\textbf{集中不等式的选择必须与被集中变量的几何结构一致}。
当理论（规范形式化）强于方法论（Hoeffding）时，集中界会因「几何不匹配」而产生系统性的紧度损失。
这种不匹配在广泛应用 SCX 于低维结构化预测问题时应予以特别关注。

\newpage

% ══════════════════════════════════════════════════════════════════════
% BIBLIOGRAPHY
% ══════════════════════════════════════════════════════════════════════
\begin{thebibliography}{99}

\bibitem{scx_theory}
SCX. \textit{Detecting Label Noise by State-Conditioned eXpertise}. SCX Theoretical System — Core Theory Volume, 2026.

\bibitem{scx_gauge}
SCX. \textit{Formal Completion of SCX Gauge Theory: $O(d)$ Lattice Gauge, Dijkgraaf-Witten Discrete TQFT, and Information-Geometric Bulk-Boundary Correspondence}. SCX Theoretical System — Gauge Field Theory Volume, v2.0, 2026.

\bibitem{hoeffding1963}
W.~Hoeffding. Probability inequalities for sums of bounded random variables. \textit{Journal of the American Statistical Association}, 58(301):13--30, 1963.

\bibitem{azuma1967}
K.~Azuma. Weighted sums of certain dependent random variables. \textit{Tohoku Mathematical Journal}, 19(3):357--367, 1967.

\bibitem{milman1971}
V.~D.~Milman. A new proof of A.~Dvoretzky's theorem on cross-sections of convex bodies. \textit{Funkcional. Anal. i Prilo\v{z}en.}, 5(4):28--37, 1971.

\bibitem{gromov1980}
M.~Gromov. Paul Levy's isoperimetric inequality. In \textit{Paul Levy's Isoperimetric Inequality}, volume~1, pages 1--25. IHES, 1980.

\bibitem{ledoux2001}
M.~Ledoux. \textit{The Concentration of Measure Phenomenon}. Mathematical Surveys and Monographs, vol.~89. American Mathematical Society, 2001.

\bibitem{boucheron2005}
S.~Boucheron, O.~Bousquet, and G.~Lugosi. Theory of classification: a survey of some recent advances. \textit{ESAIM: Probability and Statistics}, 9:323--375, 2005.

\bibitem{milman_schechtman1986}
V.~D.~Milman and G.~Schechtman. \textit{Asymptotic Theory of Finite Dimensional Normed Spaces}. Lecture Notes in Mathematics, vol.~1200. Springer, 1986.

\bibitem{anderson2003}
G.~W.~Anderson, A.~Guionnet, and O.~Zeitouni. \textit{An Introduction to Random Matrices}. Cambridge University Press, 2003.

\bibitem{mehta2004}
M.~L.~Mehta. \textit{Random Matrices}, 3rd edition. Elsevier, 2004.

\bibitem{corach2001}
G.~Corach, H.~Porta, and L.~Recht. The geometry of spaces of projections in $C^*$-algebras. \textit{Advances in Mathematics}, 101(1):59--77, 1993.

\bibitem{bump2013}
D.~Bump. \textit{Lie Groups}, 2nd edition. Graduate Texts in Mathematics, vol.~225. Springer, 2013.

\bibitem{helgason2001}
S.~Helgason. \textit{Differential Geometry, Lie Groups, and Symmetric Spaces}. Graduate Studies in Mathematics, vol.~34. American Mathematical Society, 2001.

\bibitem{jacobson1962}
N.~Jacobson. \textit{Lie Algebras}. Dover, 1962.

\bibitem{knapp2002}
A.~W.~Knapp. \textit{Lie Groups Beyond an Introduction}, 2nd edition. Progress in Mathematics, vol.~140. Birkh\"auser, 2002.

\bibitem{cheeger1975}
J.~Cheeger and D.~G.~Ebin. \textit{Comparison Theorems in Riemannian Geometry}. North-Holland, 1975.

\bibitem{gallot2004}
S.~Gallot, D.~Hulin, and J.~Lafontaine. \textit{Riemannian Geometry}, 3rd edition. Springer, 2004.

\bibitem{cornea2003}
O.~Cornea, G.~Laptev, N.~O'Shea, and C.~Viterbo. \textit{Lusternik-Schnirelmann Category}. Mathematical Surveys and Monographs, vol.~103. American Mathematical Society, 2003.

\bibitem{atkins1970}
M.~F.~Atiyah and I.~G.~Macdonald. \textit{Lusternik-Schnirelmann category and the cohomology of groups}. In \textit{Proceedings of the Conference on Algebraic Topology}, 1970.

\end{thebibliography}

\end{CJK}
\end{document}
